% ═══════════════════════════════════════════════════════════════
% KURZTEST MUSTERLÖSUNG: Reaktion von Metallen mit Sauerstoff
% Thema 2.4 | Chemie Klasse 8 | 15 Minuten | 20 BE
% ═══════════════════════════════════════════════════════════════

\documentclass[11pt, a4paper]{article}

\usepackage[utf8]{inputenc}
\usepackage[T1]{fontenc}
\usepackage[ngerman]{babel}

\usepackage{helvet}
\renewcommand{\familydefault}{\sfdefault}

\usepackage[margin=2cm, top=2.5cm, bottom=2cm]{geometry}
\usepackage{enumitem}
\usepackage{array}
\usepackage{tabularx}
\usepackage{booktabs}

\usepackage{amsmath, amssymb}

\usepackage{fancyhdr}
\usepackage{lastpage}

\usepackage{xcolor}
\usepackage{tcolorbox}
\tcbuselibrary{skins}

% ═══════════════════════════════════════════════════════════════
% FARBEN
% ═══════════════════════════════════════════════════════════════

\definecolor{chemie}{HTML}{2a7a4b}
\definecolor{hellgrau}{HTML}{F5F5F5}
\definecolor{mittelgrau}{HTML}{888888}
\definecolor{dunkelgrau}{HTML}{444444}
\definecolor{loesung}{HTML}{c62828}

% ═══════════════════════════════════════════════════════════════
% KOPFZEILE
% ═══════════════════════════════════════════════════════════════

\pagestyle{fancy}
\fancyhf{}
\renewcommand{\headrulewidth}{0.5pt}
\renewcommand{\footrulewidth}{0pt}
\fancyhead[L]{\textbf{Chemie} | Kurztest: Metalloxide | \textcolor{loesung}{\textbf{MUSTERLÖSUNG}}}
\fancyhead[R]{Klasse 8}
\fancyfoot[R]{\small\textcolor{mittelgrau}{\thepage}}

% ═══════════════════════════════════════════════════════════════
% BEFEHLE
% ═══════════════════════════════════════════════════════════════

\newcommand{\punkte}[1]{\hfill\textcolor{mittelgrau}{/#1}}
\newcommand{\lsg}[1]{\textcolor{loesung}{\textbf{#1}}}

\setlength{\parindent}{0pt}
\setlength{\parskip}{6pt}

% ═══════════════════════════════════════════════════════════════
% DOKUMENT
% ═══════════════════════════════════════════════════════════════

\begin{document}

% ═══════════════════════════════════════════════════════════════
% KOPFBEREICH
% ═══════════════════════════════════════════════════════════════

\begin{tcolorbox}[colback=loesung!10, colframe=loesung, boxrule=1pt, arc=0pt, left=8pt, right=8pt, top=6pt, bottom=6pt]
\begin{tabular}{@{}p{8cm}p{6cm}@{}}
\textbf{\large Kurztest: Reaktion von Metallen mit Sauerstoff} & \textcolor{loesung}{\textbf{MUSTERLÖSUNG}} \\[0.3em]
Thema 2.4 | 15 Minuten | 20 BE &
\end{tabular}
\end{tcolorbox}

\vspace{0.8em}

% ═══════════════════════════════════════════════════════════════
% AUFGABE 1: Begriffe (AFB I)
% ═══════════════════════════════════════════════════════════════

\textbf{Aufgabe 1:} \textit{Nenne} die passenden Fachbegriffe. \punkte{4}

\begin{enumerate}[label=\alph*), leftmargin=2em, itemsep=0.5em]
    \item Die Reaktion eines Stoffes mit Sauerstoff heißt \lsg{Oxidation}.
    \item Eine chemische Verbindung aus einem Element und Sauerstoff heißt \lsg{Oxid}.
    \item Die langsame Oxidation von Metallen im Alltag nennt man \lsg{Korrosion}.
    \item Die schützende Oxidschicht bei Aluminium nennt man \lsg{Passivierung}.
\end{enumerate}

\vspace{0.5em}

% ═══════════════════════════════════════════════════════════════
% AUFGABE 2: Wortgleichung (AFB I)
% ═══════════════════════════════════════════════════════════════

\textbf{Aufgabe 2:} \textit{Ergänze} die Wortgleichungen. \punkte{3}

\begin{enumerate}[label=\alph*), leftmargin=2em, itemsep=0.8em]
    \item Magnesium + Sauerstoff $\longrightarrow$ \lsg{Magnesiumoxid}
    \item \lsg{Eisen} + Sauerstoff $\longrightarrow$ Eisenoxid
    \item Kupfer + \lsg{Sauerstoff} $\longrightarrow$ Kupferoxid
\end{enumerate}

\vspace{0.5em}

% ═══════════════════════════════════════════════════════════════
% AUFGABE 3: Massenerhaltung (AFB II)
% ═══════════════════════════════════════════════════════════════

\textbf{Aufgabe 3:} \textit{Berechne} und \textit{erkläre}. \punkte{5}

12 g Magnesium reagieren mit 8 g Sauerstoff.

\begin{enumerate}[label=\alph*), leftmargin=2em, itemsep=0.3em]
    \item Wie viel Gramm Magnesiumoxid entstehen? \punkte{1}

    \vspace{0.3em}
    Antwort: \lsg{20} g \quad {\small\textcolor{loesung}{(12 g + 8 g = 20 g)}}

    \item Erkläre, warum das Produkt schwerer ist als das Metall allein. \punkte{4}

    \vspace{0.3em}
    \textcolor{loesung}{%
    \textbf{Erwartete Antwort:}
    \begin{itemize}[leftmargin=1.5em, itemsep=1pt]
        \item Das Metall reagiert mit Sauerstoff aus der Luft (1 BE)
        \item Die Masse des Sauerstoffs wird zum Metall hinzugefügt (1 BE)
        \item Gesetz der Massenerhaltung: Keine Atome gehen verloren (1 BE)
        \item Atome werden nur neu kombiniert/angeordnet (1 BE)
    \end{itemize}
    }
\end{enumerate}

\vspace{0.5em}

% ═══════════════════════════════════════════════════════════════
% AUFGABE 4: Korrosion (AFB II)
% ═══════════════════════════════════════════════════════════════

\textbf{Aufgabe 4:} \textit{Vergleiche} die Oxidation von Eisen und Aluminium. \punkte{4}

\begin{tabular}{|p{3cm}|p{5.5cm}|p{5.5cm}|}
\hline
& \textbf{Eisen} & \textbf{Aluminium} \\
\hline
Oxidschicht-Struktur & \lsg{porös, blättert ab} & \lsg{dicht, fest anliegend} \\[0.8em]
\hline
Schutzwirkung? & \lsg{Nein, Korrosion schreitet fort} & \lsg{Ja, Passivierung schützt} \\[0.8em]
\hline
\end{tabular}

{\small\textcolor{loesung}{Je Feld 1 BE, insgesamt 4 BE}}

\vspace{0.5em}

% ═══════════════════════════════════════════════════════════════
% AUFGABE 5: Formelgleichung (AFB II) ★★
% ═══════════════════════════════════════════════════════════════

\textbf{Aufgabe 5 (★★):} \textit{Schreibe} die Formelgleichung für die Verbrennung von Magnesium. \punkte{4}

\vspace{0.3em}
\hspace{2em} \lsg{2} Mg + \lsg{O$_2$} $\longrightarrow$ \lsg{2} \lsg{MgO}

\vspace{0.3em}
\textcolor{loesung}{%
\textbf{Bewertung:}
\begin{itemize}[leftmargin=1.5em, itemsep=1pt]
    \item Richtige Formeln (Mg, O$_2$, MgO): 2 BE
    \item Richtige Koeffizienten (2, 1, 2): 2 BE
\end{itemize}
}

\vspace{1em}

% ═══════════════════════════════════════════════════════════════
% BEWERTUNGSSCHLÜSSEL
% ═══════════════════════════════════════════════════════════════

\begin{tcolorbox}[colback=hellgrau, colframe=chemie, boxrule=1pt, arc=0pt, left=8pt, right=8pt, top=6pt, bottom=6pt]
\textbf{Bewertungsschlüssel:}

\begin{tabular}{@{}llllll@{}}
\textbf{BE} & 20--18 & 17--14 & 13--10 & 9--6 & 5--0 \\
\textbf{Note} & sehr gut (1) & gut (2) & befriedigend (3) & ausreichend (4) & mgh./ungen. (5/6)
\end{tabular}

\vspace{0.5em}
\textbf{AFB-Verteilung:} I: 7 BE (35\%) | II: 13 BE (65\%) | III: 0 BE
\end{tcolorbox}

\end{document}
